% ৯. প্রোগ্রামের সংগঠন
\section{প্রোগ্রামের সংগঠন}

\begin{learningbox}
এই টপিক শেষে শিক্ষার্থী একটি C program-এর প্রধান অংশ এবং IPO model ব্যাখ্যা করতে পারবে।
\end{learningbox}

\begin{definitionbox}{Program organization}
একটি program-এর বিভিন্ন অংশ যেমন comment, header, declaration, main function, input, processing, output ও return statement সুশৃঙ্খলভাবে সাজানোর পদ্ধতিকে program organization বলা যায়।
\end{definitionbox}

\begin{center}
\fbox{\parbox{0.88\textwidth}{
\centering
\textbf{IPO model placeholder}\\[4pt]
Input $\rightarrow$ Processing $\rightarrow$ Output
}}
\end{center}

\begin{keypointbox}{C program-এর সাধারণ অংশ}
\begin{enumerate}
\item Documentation/comment
\item Preprocessor directive
\item Global declaration
\item \texttt{main()} function
\item Variable declaration
\item Input, processing, output
\item \texttt{return 0;}
\end{enumerate}
\end{keypointbox}

\begin{boardbox}{বোর্ড ফোকাস}
C program structure আঁকতে \texttt{\#include <stdio.h>} এবং \texttt{int main()} অংশ সাধারণত লিখতে হয়।
\end{boardbox}
